\documentclass[11pt,a4paper]{article}
\usepackage[T1]{fontenc}
\usepackage[utf8]{inputenc}
\usepackage{lmodern,amsmath,amssymb}
\usepackage[margin=25mm]{geometry}
\usepackage[hidelinks]{hyperref}
\hypersetup{pdftitle={Shared resources and exclusive receiver events},pdfauthor={navstokgap research working notes}}
\newcommand{\tightlist}{\setlength{\itemsep}{0pt}\setlength{\parskip}{0pt}}
\setlength{\emergencystretch}{3em}
\setcounter{secnumdepth}{0}
\begin{document}
\hypertarget{a-shared-release-produces-exclusive-fringe-weighted-events}{%
\section{A shared release produces exclusive fringe-weighted
events}\label{a-shared-release-produces-exclusive-fringe-weighted-events}}

An autonomous classical jump receiver can produce exactly one eventual
event with probabilities equal to normalized interference energies. It
does so by admitting one shared release and prescribing rates linear in
those energies. Finite observation time reveals a signal-dependent
efficiency; spatially separated local releases require a communication
or common-control mechanism. Neither construction selects a universal
action constant.

Status: Q11 exploratory model and written derivation. This is an
effective stochastic receiver, not a microscopic Hamiltonian derivation.
Its discrete release states, stored energy and rate law are explicit
model inputs.

\hypertarget{autonomous-competition-and-energy-accounting}{%
\subsection{1. Autonomous competition and energy
accounting}\label{autonomous-competition-and-energy-accounting}}

First load the two classical work records from Q08 into ideal storage:

\[e_++e_-=E>0,\qquad e_\pm=E(1\pm z)/2,\quad -1\le z\le1.\]

The following receiver stage has constant stored inputs. Let R be a
ready state and \(C_+,C_-\) be two absorbing record states. Choose a
transduction constant \(\alpha>0\) of units
\((\text{energy}\,\text{time})^{-1}\) and transitions

\[R\xrightarrow{r_+=\alpha e_+}C_+,\qquad
R\xrightarrow{r_-=\alpha e_-}C_-. \tag{1}\]

This time-homogeneous Markov model is autonomous after preparation. Its
absorbing architecture stipulates that the first channel consumes the
single permission to release; no second transition remains. The stored
signal energies act as kinetic controls and are not automatically the
emitted record energy.

For explicit energy bookkeeping, put one releasable amount \(\Delta>0\)
in a central reservoir. Either transition lowers that reservoir by
\(\Delta\) and increases the selected output's energy plus heat by
\(\Delta\). Signal storage is unchanged in this idealization. Resetting
the released reservoir requires replacement energy \(\Delta\), or
retrieval from an explicitly included output store. The architecture
supplies single use; energy conservation alone does not imply the state
graph (1). Permanent storage, reset efficiency and a physical
implementation of the rate control remain additional premises.

\hypertarget{exact-probabilities-and-their-limit}{%
\subsection{2. Exact probabilities and their
limit}\label{exact-probabilities-and-their-limit}}

Write \(\Gamma=r_++r_-=\alpha E\). The survival equation is
\(\dot p_R=-\Gamma p_R\), \(p_R(0)=1\). Integrating each exit flux gives

\[p_R(t)=e^{-\Gamma t},\qquad
p_\pm(t)=\frac{r_\pm}{\Gamma}(1-e^{-\Gamma t}). \tag{2}\]

Double records have probability zero. Conditional on one event by any
fixed \(t>0\), or unconditionally after waiting indefinitely,

\[\Pr(C_\pm\mid\text{event})=\frac{e_\pm}{E}=\frac{1\pm z}{2}. \tag{3}\]

These fringe weights follow from the assumed linear hazards and
competition; no quantum probability rule was used to obtain them. They
also do not establish quantum state structure. Replacing the hazards by
\(\alpha_p e_\pm^p\), with the corresponding units and \(p>0\), gives
weights \(e_\pm^p/(e_+^p+e_-^p)\). Exclusivity therefore does not select
linearity.

If the incident displacement is attenuated by \(\eta>0\), then \(e_\pm\)
scale by \(\eta^2\) and z is unchanged. Equation (3) persists, but at a
fixed deadline T the efficiency becomes \(1-\exp(-\alpha\eta^2ET)\) and
tends to zero. Requiring efficiency at least \(1-\varepsilon\),
\(0<\varepsilon<1\), gives

\[E\ge\frac{\log(1/\varepsilon)}{\alpha T}. \tag{4}\]

At fixed carrier frequency \(\omega\) in Q08's narrow-band regime, the
associated wave-action bound is approximately
\(\log(1/\varepsilon)/(\alpha T\omega)\). It retains receiver coupling,
deadline, frequency and error tolerance. The indefinite-wait and
zero-signal limits do not commute. At E=0 there is no event.

\hypertarget{what-spatial-separation-changes}{%
\subsection{3. What spatial separation
changes}\label{what-spatial-separation-changes}}

Equation (1) is consistent as a co-located central arbitration model:
both stored inputs are available before the release. Sending the output
to a remote record then takes finite travel time. It is not an
instantaneous inhibition law for distant detectors that already register
events independently.

To test that alternative, let the two sites initially have independent
Poisson trigger clocks with rates \(r_+,r_-\). After the first event, an
inhibition signal reaches the other site after a fixed delay
\(\delta>0\). Each site can fire only once. On an observation interval
long enough to include that delay, memorylessness gives the exact
eventual double-event probability

\[P_2=\frac{r_+}{\Gamma}(1-e^{-r_-\delta})+
       \frac{r_-}{\Gamma}(1-e^{-r_+\delta}). \tag{5}\]

It is positive whenever both rates and the delay are positive. If their
separation is d and inhibition speed is at most c, \(\delta\ge d/c\) in
this specified architecture. In the balanced case \(r_+=r_-=\Gamma/2\),

\[P_2=1-e^{-\Gamma\delta/2}. \tag{6}\]

Combining a first-event deadline T with \(P_2\le\varepsilon_2\) and
\(\Pr(\text{no first event by }T)\le\varepsilon_0\) requires

\[\frac{\log(1/\varepsilon_0)}{T}\le\Gamma
\le\frac{-2\log(1-\varepsilon_2)}{\delta}. \tag{7}\]

The double-event criterion here counts events during the full inhibition
window, even when it extends past T. Short-gate censoring would change
that criterion. Equation (7) is a timing compatibility condition, not an
action gap or a bound for every classical architecture. A central
arbiter, prior correlated readiness, or a different propagation law
changes a premise.

\hypertarget{decision}{%
\subsection{4. Decision}\label{decision}}

Shared-resource competition successfully turns classical fringe energies
into exclusive conditional event probabilities in a declared model. The
resource preparation and linear kinetic response do substantive work.
Replacing central arbitration by local firing with delayed inhibition
adds a calculable efficiency/coincidence trade-off. Both mechanisms
retain adjustable energy and clock scales.

After Q09--Q11, park further detector tuning. The decisive next
construction is a pair of independently prepared local receivers with no
communication during the gate: ask whether independent monotone local
response to a common classical pulse energy permits suppressed
coincidences without selection of gates. This changes the
common-resource premise rather than refining another rate constant. Keep
source conditioning and detector independence explicit; any resulting
exclusion must be limited to that stated class.

\hypertarget{source-boundary}{%
\subsection{Source boundary}\label{source-boundary}}

D. T. Gillespie, \emph{Exact stochastic simulation of coupled chemical
reactions}, Journal of Physical Chemistry \textbf{81}, 2340--2361
(1977), \href{https://doi.org/10.1021/j100540a008}{DOI}, is the primary
discovery lead for the standard competing-hazard construction. Coverage
here is one query and publisher metadata, not a full-text audit or
simulation. Equations (2)--(7) are obtained directly by integrating
exponential survival probabilities. Independent written review and
bounded librarian comparison are required before accepted-ledger
promotion. No numerical or symbolic scripts were used.
\end{document}
